%=====================================================================
%  Chapter 2 — Concepts on Bipedal Walking
%=====================================================================
\chapter{Concepts on Bipedal Walking}
\label{chap:walking}

\section{Introduction}
Bipedal walking, or walking on two legs, is a complex form of locomotion
that involves the coordination of multiple joints to achieve stability,
balance, and forward motion. Understanding its principles is essential
for designing and controlling a humanoid robot. This chapter introduces
the fundamental concepts of bipedal walking: the definition of walking,
the motion planes, the centre of mass, the Zero-Moment-Point, the
support polygon, the stability criteria, the types of walking, and the
gait cycle that together underpin the locomotion work presented later in
this report.

\section{Walking Definition}
Walking can be defined as the action of moving by successive movements
and support of the legs and feet without ever leaving the ground
entirely. This definition does not constrain the number of limbs, so
walking may be bipedal, quadrupedal, and so on; in this work the term
refers to \emph{bipedal} walking, since the objective is to reproduce
human locomotion. Walking is therefore a cyclic phenomenon: each foot
alternately swings through the air and then serves as support, either
alone or together with the other foot. This cyclic nature is formalised
through the notion of gait \emph{cycles} and \emph{phases}, developed in
the following sections.

\section{Motion Planes}
The motion planes describe the movements associated with walking, as
shown in Figure~\ref{fig:planes}. The \textbf{sagittal} plane
corresponds to a side view of the walker, the \textbf{horizontal} plane
to a top view, and the \textbf{frontal} plane to a front view. The
walker moves primarily in the sagittal plane, but movements such as hip
sway, lateral displacement, and pelvic rotation occur in the other two
planes and enhance the fluidity of walking.

\begin{figure}[ht]
  \centering
  \includegraphics[width=0.75\textwidth]{fig_motion_planes.png}
  \caption{Motion planes: sagittal, horizontal, and frontal.}
  \label{fig:planes}
\end{figure}

\section{Robot Control and Stability}
Bipedal walking relies on several key concepts that ensure stability and
efficiency. These concepts are essential for understanding how to design
and control a bipedal robot.

\subsection{Centre of Mass (CoM)}
The centre of mass (CoM), also known as the centre of gravity, is the
point at which the mass of a body or system of bodies can be considered
to be concentrated. For a rigid body, it is the point at which the body
can be balanced without rotation under the influence of gravity alone.
Mathematically, it is the weighted average of the positions of all the
mass elements of the system. The CoM plays a central role in the study
of the mechanics, dynamics, and stability of a walking robot.

\subsection{Zero-Moment-Point (ZMP)}
The Zero-Moment-Point (ZMP) is a criterion used to analyse the stability
of a bipedal or multi-legged robot during dynamic motion. The ZMP is the
point on the ground where the total moment of the forces acting on the
robot's feet has no horizontal component --- that is, the point where the
ground-reaction forces and the inertial forces acting on the robot's body
balance each other out.

Consider the support foot during the single-support phase
(Figure~\ref{fig:zmp}). The effect of the parts above the support foot
can be replaced by a force $\mathbf{F}$ and a moment $\mathbf{M_A}$ at
the ankle, while the ground reaction $\mathbf{R}$ and moment
$\mathbf{M_S}$ balance the system at the contact point $S$. Because the
ground can only act on the foot through friction and a normal force, the
horizontal moments vanish at this point --- hence the term
\emph{Zero-Moment-Point}. The ZMP is thus the point under the foot where
the ground reaction balances the horizontal moments applied at the
ankle. To maintain dynamic balance, the ZMP must remain inside the
support polygon; if the support area is too small to compensate the
horizontal moment, the robot loses dynamic balance and the ZMP is no
longer defined (a fictitious ZMP may then be computed for analysis).

\begin{figure}[ht]
  \centering
  \includegraphics[width=0.85\textwidth]{fig_zmp.png}
  \caption{General case of the ZMP: (a) three-dimensional case at the
           ankle, (b) planar case where the moment $M_{Ax}$ is balanced
           by shifting the point of action of the reaction force $R_z$.}
  \label{fig:zmp}
\end{figure}

\subsection{Support Polygon}
The support polygon is one of the stability criteria for humanoid
walking. It is defined by the outer contour of the points of contact
between the robot's feet and the ground. As illustrated in
Figure~\ref{fig:support}, it takes two forms during the walking cycle:
the \textbf{single-support} polygon, reduced to the contour of a single
foot, and the larger \textbf{double-support} polygon, spanning both feet
when they are simultaneously in contact with the ground.

\begin{figure}[ht]
  \centering
  \includegraphics[width=0.55\textwidth]{fig_support_polygon.png}
  \caption{Support polygon in the single- and double-support phases.}
  \label{fig:support}
\end{figure}

\subsection{Stability Criteria}
Two complementary aspects of stability are considered:
\begin{itemize}[itemsep=2pt,topsep=4pt]
  \item \textbf{Posture stability} --- the ability to maintain a stable
        standing position with minimal jerk. Joint positions and the
        centre of gravity are the main variables involved.
  \item \textbf{Gait stability} --- the ability to maintain balance
        during locomotion. The key variables here are step length,
        foothold positions, and ZMP data.
\end{itemize}

\section{Types of Bipedal Walking}
Walking can be categorised into two types:
\begin{itemize}[itemsep=2pt,topsep=4pt]
  \item \textbf{Static walking.} The posture is stable at all times: the
        body weight is always supported by at least one foot, and the
        projection of the CoM on the ground always remains within the
        support polygon.
  \item \textbf{Dynamic walking.} A succession of stable and unstable
        postures. Instability is deliberately created forward when the
        projection of the CoM moves outside the support polygon. Dynamic
        walking is the natural human gait; the resulting instability is
        corrected by the next step.
\end{itemize}

\section{Robot Gait Generation}

\subsection{Walking-Cycle Components}
Ignoring the very first and very last steps to and from a standing
posture, the walking cycle is composed of two alternating phases:
\begin{description}[itemsep=2pt,topsep=4pt]
  \item[Double-Support Phase (DSP)] The robot is in contact with the
        floor through both legs. This phase lets the robot shift its ZMP
        into a favourable position for the next step without losing
        balance; it is therefore called the \emph{stabilisation} phase.
  \item[Single-Support Phase (SSP)] Also known as the swing phase, the
        robot is in contact with the floor through a single leg. The leg
        in contact is the \emph{support} leg and the other is the
        \emph{swinging} leg.
\end{description}

\subsection{Walking-Cycle Phases}
Within these components, the cycle progresses through the following
phases:
\begin{itemize}[itemsep=2pt,topsep=4pt]
  \item \textbf{Impact / contact phase (DSP)} --- the heel of the front
        foot comes into contact with the floor.
  \item \textbf{Lift-off / swing phase (SSP)} --- during the transition
        from double to single support, the swinging foot is lifted off
        the ground; the robot transfers its weight onto the support foot
        and starts swinging the rear leg forward.
  \item \textbf{Catch phase (DSP)} --- the swinging foot comes back into
        contact with the floor, returning the robot to double support.
\end{itemize}
The roles of the two legs are then swapped and the cycle repeats.

\subsection{Gait Engine}
The gait engine, also referred to as the walking-pattern generator, is
the software component responsible for producing a sequence of feasible
joint positions that reproduce a walking motion. It accounts for the
robot's physical limitations, its preferred walking speed, and the
desired gait style to compute the best sequence of joint angles and body
motions for stable and efficient walking. Gait engines fall into two
main families:
\begin{itemize}[itemsep=2pt,topsep=4pt]
  \item \textbf{ZMP-based approach} --- requires comprehensive knowledge
        of the robot's parameters and uses the ZMP as a feedback signal
        to adjust joint angles and torques.
  \item \textbf{Inverted-pendulum approach} --- operates with limited
        knowledge of the robot model and uses a closed-loop feedback
        based on the inverted-pendulum abstraction.
\end{itemize}
The gait generator developed in this work follows the ZMP-based approach
and is detailed in Chapter~\ref{chap:kinematics}.

\section{Walking Primitive}
A walking primitive is a fundamental building block that generates the
motion of a robot's gait. It defines characteristics such as step
length, step height, and step frequency, taking into account the robot's
physical constraints, its stability requirements, and the desired
locomotion behaviour.

\section{Conclusion}
This chapter introduced the fundamental concepts of bipedal walking,
providing the vocabulary and the stability criteria --- centre of mass,
Zero-Moment-Point, and support polygon --- that govern stable, human-like
locomotion, together with the structure of the gait cycle. These
concepts form the theoretical foundation for the hardware, firmware, and
control chapters that follow. The next chapter presents the hardware and
mechatronic design of the Abbes robot.
